\label{...}

\import{chapters/2. Base-Knowledge/2.1 Recommender-System-Big-Data/}{2.1.1 Recommender-Systems.tex}
\import{chapters/2. Base-Knowledge/2.1 Recommender-System-Big-Data/}{2.1.2 Session-based-Recommendation.tex}

% \section{Kiến thức về Dữ liệu lớn}
% \import{chapters/2. Base-Knowledge/2.1 Recommender-System-Big-Data/}{2.1.3 Data-Pipeline.tex}

% \section{Kiến thức về Container \& Orchestration}
% \label{...}

% \import{chapters/2. Base-Knowledge/2.2 Container-Orchestration/}{2.2.1 Containerization.tex}
% \import{chapters/2. Base-Knowledge/2.2 Container-Orchestration/}{2.2.2 Orchestration.tex}
% \import{chapters/2. Base-Knowledge/2.2 Container-Orchestration/}{2.2.3 Infrastructure-as-Code.tex}


% \section{Các công nghệ được sử dụng}
% \label{...}

% \import{chapters/2. Base-Knowledge/2.3 Technologies/}{2.3.1 Florus.tex}
% \import{chapters/2. Base-Knowledge/2.3 Technologies/}{2.3.2 Docker.tex}
% \import{chapters/2. Base-Knowledge/2.3 Technologies/}{2.3.3 Kubernetes-Minikube.tex}
% \import{chapters/2. Base-Knowledge/2.3 Technologies/}{2.3.4 Helm.tex}
% \import{chapters/2. Base-Knowledge/2.3 Technologies/}{2.3.5 DB-name.tex}


\section{Kết chương}

Tóm lại, hai yếu tố chính định hình thiết kế của bài tập là: giới hạn thực tế của ngôn ngữ (đặc biệt là đệ quy) và mục tiêu xây dựng một môi trường độc lập, có thể mở rộng cho AI sau này.

Hệ thống được thiết kế theo hướng tách rõ logic và giao diện, đồng thời ưu tiên kiểm soát bộ nhớ và độ ổn định khi xử lý flood fill. Những yêu cầu này dẫn đến việc cần đánh giá lại các thuật toán duyệt đồ thị phổ biến, nhằm chọn ra phương án phù hợp nhất. Nội dung đó sẽ được phân tích trong chương tiếp theo.